\metatitle{Colored lattice-path matroids}
\metadescription{Colored lattice-path matroids defined by prefix-count bounds,
their assignment polytopes, and their cubical regions in the quota box.}

\section[coloredLatticePaths]{Colored lattice-path matroids}

An ordinary lattice-path matroid consists of the north-step sets of all paths
between two fixed north--east paths \cite{BoninMierNoy2003}. The two boundary
paths give lower and upper bounds on the number of north steps in every
prefix. We use the same description with one prefix count for each color.

Fix quotas $r_1,\dotsc,r_k$ with sum $n$. A word
$w=w_1\dotsm w_n$ with this content determines a path in $\mathbb Z^k$ by
setting
\[
 \defin{h_c(w,j)}\coloneqq |\{i\leq j:w_i=c\}|,
 \qquad
 \defin{p_w(j)}\coloneqq
 \bigl(h_1(w,j),\dotsc,h_k(w,j)\bigr).
\]
Thus $p_w(0)=0$, $p_w(n)=(r_1,\dotsc,r_k)$, and the $j$th step is the
coordinate vector indexed by $w_j$.

Choose a set $L_j\subseteq[k]$ of allowed colors at each position $j$.
At selected cuts $j\in[n-1]$, choose integral prefix bounds
\[
 \ell_c(j)\leq h_c(w,j)\leq u_c(j)
 \qquad(c\in[k]).
\]
We take every word with the prescribed content that respects the allowed
colors and all the prefix bounds. A color \defin{crosses the cut $j$} if it
is allowed at some position on each side of the cut.

\begin{definition}[Colored lattice-path matroid]
A \defin{colored lattice-path matroid} is a nonempty family obtained from
allowed colors and prefix bounds as above, where at most two colors cross
each cut carrying a prefix bound.
\end{definition}

Different pairs of colors may cross different cuts. The definition therefore
allows a path to move in all $k$ coordinate directions, although each
boundary condition concerns a part of the path in which at most two colors
remain mobile.

\begin{theorem}[Two-color cuts]
Every colored lattice-path matroid is a colored matroid.
\end{theorem}

\begin{proof}
We induct on the number of constrained cuts. With no constrained cuts, the
family is an allowed-color family, so symmetric cycle exchange follows from
the balanced disagreement graph.

Let $j$ be the first constrained cut, and compare two words $u,v$. A color
that does not cross $j$ has the same number of occurrences in $[j]$ in both
words. If the two prefix-count vectors differ, their difference is therefore
supported on the two colors crossing $j$, and has the form $(d,-d)$.

If the prefix counts agree, take $d=0$ and add no positions. Otherwise, name
the crossing colors $a,b$ so that $u$ has $d$ more occurrences of $a$ on the
left and $v$ has $d$ more occurrences of $b$. Append $b^d$ to the left part
of $u$ and $a^d$ to the left part of $v$, and prepend the opposite
assignments to their right parts. The two extended words now have equal color
quotas on each side.
The left side has no constrained cut. On the right, if $r_c'$ is the extended
quota of color $c$, the bound at a later cut $h$ is translated using
\[
 h_c(w,h)=r_c-r_c'+h_c(w',d+h-j),
\]
where $w'$ is the extended right part. The temporary colors already crossed
any later cut at which they remain mobile: each was allowed on the original
left side of that cut. Thus the translation introduces no new crossing color,
and we may apply symmetric cycle exchange on each side by induction.

If the exchange cycle through the prescribed disagreement avoids the
temporary positions, it is already an exchange in the original words.
Otherwise it uses one temporary arrow. The exchange cycle through the
opposite temporary arrow on the other side joins to it after the two
temporary arrows are removed. Since only the two crossing colors occur on
both sides, the resulting directed cycle is simple. The two offspring satisfy
the original allowed-color conditions and quotas. At the cut $j$, their two
crossing-color counts lie between the corresponding counts of $u$ and $v$;
the translated bounds handle all later cuts.
\end{proof}

\begin{proposition}[Two colors]
Suppose there are two colors and both are allowed at every position. Then the
colored lattice-path matroids are precisely ordinary lattice-path matroids,
with a basis recorded by the positions of one color.
\end{proposition}

\begin{proof}
Every cut is crossed by at most two colors. The lower and upper prefix counts
of color $1$ are the height functions of the lower and upper boundary paths.
Conversely, the paths between two such boundaries are exactly the binary
words satisfying these prefix bounds.
\end{proof}

\begin{example}[Different pairs at different cuts]
Take three colors with quotas $(2,2,2)$. Allow colors $1,2$ in the first two
positions, all three colors in positions $3,4$, and colors $2,3$ in the last
two positions. At the cut after position $2$, impose
\[
 h_1(w,2)\geq1,
\]
and at the cut after position $4$, impose
\[
 h_3(w,4)\leq1.
\]
Only colors $1,2$ cross the first cut, while only colors $2,3$ cross the
second. The word $\mathtt{123123}$ is one of the resulting paths.
\end{example}

The restriction on constrained cuts is needed. With quotas $(2,1,1)$, take
all words satisfying
\[
 h_1(w,2)\leq1.
\]
This removes the paths passing through $(2,0,0)$. The words
$\mathtt{1213}$ and $\mathtt{3121}$ are allowed. The only cyclic exchange
through their first position switches positions $1$ and $4$, and one of the
two offspring is $\mathtt{1123}$. Hence this ten-word path family is not a
colored matroid.

\subsection[pathPolytope]{The path polytope}

For every position $j$ and color $c$, let $x_{j,c}$ record whether the $j$th
step has color $c$. The assignment polytope of a colored lattice-path matroid
has the following description:
\[
\begin{aligned}
 x_{j,c}&\geq0,
 &x_{j,c}&=0 &&\text{if }c\notin L_j,\\
 \sum_{c=1}^k x_{j,c}&=1,
 &\sum_{j=1}^n x_{j,c}&=r_c,\\
 \ell_c(s)&\leq\sum_{j\leq s}x_{j,c}\leq u_c(s)
 &&&\text{at every constrained cut }s.
\end{aligned}
\]

\begin{proposition}
The polytope above is integral. Its vertices are the words of the colored
lattice-path matroid.
\end{proposition}

\begin{proof}
Use one node of supply one for each position. For every allowed pair $(j,c)$,
send its flow to time $j$ on a chain belonging to color $c$. The flow on the
chain arc from time $s$ to time $s+1$ is the color-$c$ prefix count through
position $s$, so the prefix bounds become lower and upper arc capacities.
The last node of the color-$c$ chain has demand $r_c$.

This is a network-flow polytope with integral supplies, demands, and
capacities, and is therefore integral. Its integral flows have one outgoing
color at each position and are exactly the admissible words.
\end{proof}

\subsection[pathRegions]{Cubical regions}

Let
\[
 Q=[0,r_1]\times\dotsm\times[0,r_k].
\]
At distance $s$ from the origin, a path lies in the section
\[
 \defin{\Delta_s}\coloneqq\left\{a\in\mathbb Z_{\geq0}^k:
 a_1+\dotsb+a_k=s,\quad a_c\leq r_c\right\}.
\]
For a nonempty prefix-bound family $\mathcal F$, let
\[
 \ell_c(s)=\min_{w\in\mathcal F}h_c(w,s),
 \qquad
 u_c(s)=\max_{w\in\mathcal F}h_c(w,s).
\]
Replacing the original bounds by these actual extrema does not change
$\mathcal F$. Its canonical vertex region is
\[
 R_{\mathcal F}=\left\{a\in Q\cap\mathbb Z^k:
 \begin{array}{c}
  \ell_c(|a|)\leq a_c\leq u_c(|a|)\\
  \text{for every }c
 \end{array}\right\}.
\]
Here $|a|=a_1+\dotsb+a_k$.

For each color $c$, collapse all the other directions to one direction:
\[
 \defin{\pi_c(a)}=\left(a_c,\sum_{d\neq c}a_d\right).
\]
This maps $Q$ onto the rectangle
$[0,r_c]\times[0,n-r_c]$. The inequalities
\[
 \ell_c(x+y)\leq x\leq u_c(x+y)
\]
cut out an ordinary planar lattice-path corridor $D_c$ in this rectangle.
Indeed, both $\ell_c$ and $u_c$ start at $0$, end at $r_c$, and increase by
either $0$ or $1$ at each step.

\begin{theorem}[Cubical-region characterization]
Every nonempty prefix-bound family has a canonical region of the form
\[
 R=Q\cap\mathbb Z^k\cap\bigcap_{c=1}^k\pi_c^{-1}(D_c),
\]
where $D_c$ is an ordinary lattice-path corridor from $(0,0)$ to
$(r_c,n-r_c)$. Conversely, any such intersection that contains a complete
permitted path defines a nonempty prefix-bound family. Equivalently, every
rank section is a simplex cut by a coordinate box:
\[
 R\cap\Delta_s=
 \left\{a\in\Delta_s:
 \begin{array}{c}
  \ell_c(s)\leq a_c\leq u_c(s)\\
  \text{for every }c
 \end{array}\right\}.
\]
\end{theorem}

\begin{proof}
Since $|a|=a_c+\sum_{d\neq c}a_d$, membership of $\pi_c(a)$ in $D_c$ is
exactly the pair of inequalities
$\ell_c(|a|)\leq a_c\leq u_c(|a|)$. Intersecting these conditions over all
colors gives the displayed intersection, and restricting to $|a|=s$ gives
the section formula. Conversely, the two boundary paths of each $D_c$ supply
the corresponding lower and upper prefix bounds.
\end{proof}

The corridors describe the vertices. The allowed-color sets describe the
edges: at rank $j-1$, retain the edge
\[
 a\longrightarrow a+e_c
\]
precisely when both endpoints lie in $R$ and $c\in L_j$. The admissible words
are exactly the monotone paths from $0$ to $(r_1,\dotsc,r_k)$ in this directed
region. Its path-filled part is obtained by retaining the vertices and edges
that occur on a complete path, and then filling every unit cube supported by
this one-skeleton.

There is also a direct description in terms of deleted orthants. Write
$a_{-c}$ for the vector obtained by omitting $a_c$, and put
$|a_{-c}|=\sum_{d\neq c}a_d$. A forbidden quadrant of the $c$th planar
corridor pulls back to a set of one of the two forms
\[
 \begin{gathered}
  \left\{a:a_c\geq A,\ |a_{-c}|\leq B\right\},\\
  \left\{a:a_c\leq A,\ |a_{-c}|\geq B\right\}.
 \end{gathered}
\]
For example, inside $Q$,
\[
 \begin{aligned}
  &\left\{a:a_c\geq A,\ |a_{-c}|\leq B\right\}\\
  &\quad=
  \bigcup_{\substack{q\in\mathbb Z_{\geq0}^{k-1}\\|q|=B}}
  \left\{a:
  \begin{array}{c}
   a_c\geq A,\\[-2pt]
   a_{-c}\leq q
  \end{array}\right\},
 \end{aligned}
\]
where $a_{-c}\leq q$ means coordinatewise inequality. The second form has
the analogous expression with all inequalities reversed. Thus a planar
corner becomes a complete family of mixed orthants:
all anchors with the same total in the other coordinates occur together.
This is the exact higher-dimensional analogue of removing everything on one
side of a corner in a skew shape.

More generally, for $S\subseteq[k]$ one may use
\[
 \pi_S(a)=\left(\sum_{c\in S}a_c,
                    \sum_{c\notin S}a_c\right).
\]
The pullbacks of planar quadrants then have one sign on the coordinates in
$S$ and the opposite sign on the coordinates outside $S$. The present
prefix-count definition is the case $S=\{c\}$.

Finally, suppose a prefix bound is imposed at the cut $s$. Every color that
does not cross this cut has coordinate either $0$ or $r_c$ throughout the
admissible section. Hence the mobile part of the section uses only the
crossing colors. For a colored lattice-path matroid there are at most two of
them, so each constrained section is an interval in a two-coordinate plane.
The active pair may change from one cut to the next.
